% =============================================================================
% figures/fig_literature_positioning.tex
%
% 2x2 literature-positioning matrix: colours and markers match the PDF version
% generated by make_additional_figures.py::make_literature_positioning().
%
% Quadrant palette (HTML codes, 7% fill opacity over white):
%   BL theoretical/static  : #4878CF (blue)
%   BR empirical/static    : #4da06d (green)
%   TL theoretical/temporal: #aaaaaa (grey, 3% opacity, no prior work)
%   TR empirical/temporal  : #E69F00 (amber)
%
% Required TikZ libraries (preamble.tex):
%   arrows.meta, shapes.geometric, calc
% =============================================================================

% Colour definitions
\definecolor{CTheory}{HTML}{4878CF}   % blue
\definecolor{CEmpSt} {HTML}{4da06d}   % green
\definecolor{CEmpTs} {HTML}{E69F00}   % amber
\definecolor{CThesis}{HTML}{3c7fd0}   % AG blue (this thesis)

\begin{tikzpicture}[>= Stealth, x=1cm, y=1cm]

% ── Layout ────────────────────────────────────────────────────────────────────
% Content box : (1.4, 0.9) to (12.6, 10.0)
% Divider     : x = 7.0,  y = 5.45
% x-axis arrow: y = 0.55
% y-axis arrow: x = 0.90

% ── Quadrant fills — diagonal gradients converging toward centre (7.0, 5.45) ──
% BL: lighter at outer corner (1.4,0.9), darker toward centre (angle=45°)
\shade[shading=axis, shading angle=45,
       left color=CTheory!4, right color=CTheory!24]
  (1.4, 0.9) rectangle (7.0, 5.45);
% BR: lighter at outer corner (12.6,0.9), darker toward centre (angle=135°)
\shade[shading=axis, shading angle=135,
       left color=CEmpSt!4, right color=CEmpSt!24]
  (7.0, 0.9) rectangle (12.6, 5.45);
% TL: very subtle, no prior work; lighter at outer corner, faint toward centre
\shade[shading=axis, shading angle=315,
       left color=black!1, right color=black!7]
  (1.4, 5.45) rectangle (7.0, 10.0);
% TR: lighter at outer corner (12.6,10.0), darker toward centre (angle=225°)
\shade[shading=axis, shading angle=225,
       left color=CEmpTs!4, right color=CEmpTs!24]
  (7.0, 5.45) rectangle (12.6, 10.0);

% ── Dividers and border ───────────────────────────────────────────────────────
\draw[dashed, gray!78, thin] (1.4, 5.45) -- (12.6, 5.45);
\draw[dashed, gray!78, thin] (7.0, 0.9)  -- (7.0, 10.0);
\draw[gray!68, thin] (1.4, 0.9) rectangle (12.6, 10.0);

% ── Axis arrows ───────────────────────────────────────────────────────────────
\draw[->, black, thin] (1.10, 0.55) -- (12.88, 0.55);
\draw[->, black, thin] (0.90, 0.72) -- (0.90, 10.18);

% ── X-axis labels ─────────────────────────────────────────────────────────────
\node[font=\small\itshape, text=black, anchor=north west] at (1.4, 0.43)
  {Theoretical};
\node[font=\small\itshape, text=black, anchor=north east] at (12.6, 0.43)
  {Empirical, AutoML-bounded};
\node[font=\small\bfseries, text=black, anchor=north] at (7.0, 0.08)
  {Rashomon set operationalisation};

% ── Y-axis labels — centred in their respective halves ────────────────────────
\node[font=\small\itshape, text=black, rotate=90, anchor=center]
  at (0.66, 3.18) {Static tabular};
\node[font=\small\itshape, text=black, rotate=90, anchor=center]
  at (0.66, 7.73) {Temporal, AutoML};
\node[font=\small\bfseries, text=black, rotate=90, anchor=south] at (0.18, 5.45)
  {Data setting};

% ── Quadrant corner labels ────────────────────────────────────────────────────
\node[font=\footnotesize\itshape, text=black,
      anchor=north west, align=left] at (1.55, 5.38)
  {Theoretical Rashomon set\\Static tabular data};
\node[font=\footnotesize\itshape, text=black,
      anchor=north west, align=left] at (7.15, 5.38)
  {Empirical Rashomon set\\Static tabular data};
\node[font=\footnotesize\itshape, text=black,
      anchor=south west, align=left] at (7.15, 5.52)
  {Empirical, AutoML\\Temporal data};

% ── Literature markers ────────────────────────────────────────────────────────

% Bottom-left: theoretical Rashomon set, static tabular, blue circles
% Clustered together (gap 0.65)
\filldraw[fill=CTheory, draw=white, line width=0.5pt]
  (2.40, 3.85) circle (3.0pt);
\node[font=\small, text=black, anchor=west] at (2.72, 3.85)
  {Fisher et al.\ (2019)};

\filldraw[fill=CTheory, draw=white, line width=0.5pt]
  (2.40, 3.20) circle (3.0pt);
\node[font=\small, text=black, anchor=west] at (2.72, 3.20)
  {Laberge et al.\ (2023)};

\filldraw[fill=CTheory, draw=white, line width=0.5pt]
  (2.40, 2.55) circle (3.0pt);
\node[font=\small, text=black, anchor=west] at (2.72, 2.55)
  {Donnelly et al.\ (2023)};

% Bottom-right: empirical, static tabular, green square
\filldraw[fill=CEmpSt, draw=white, line width=0.5pt]
  (8.50-0.12, 3.20-0.12) rectangle (8.50+0.12, 3.20+0.12);
\node[font=\small, text=black, anchor=west] at (8.82, 3.20)
  {M\"{u}ller et al.\ (2023)};

% Top-right: Zhao and Ma (2025), amber diamond, italic label
\node[diamond, fill=CEmpTs, draw=white, minimum size=9pt, inner sep=0pt]
  at (8.20, 6.90) {};
\node[font=\small\itshape, text=black, anchor=west] at (8.55, 6.90)
  {Zhao \& Ma (2025)$^{\dagger}$};

% Top-right: This thesis, star inside a drawn box
\node[star, star points=5, star point ratio=2.25,
      fill=CThesis, draw=white, minimum size=14pt, inner sep=0pt]
  at (9.20, 8.40) {};
\node[draw=CThesis, fill=CThesis!12, rounded corners=3pt,
      inner sep=5pt, font=\small\bfseries, text=CThesis,
      anchor=west] at (9.58, 8.40)
  {This thesis};

% ── Footnote ──────────────────────────────────────────────────────────────────
\node[font=\tiny\itshape, text=black,
      anchor=north west, text width=11.0cm, align=left]
  at (1.40, -0.30)
  {$^{\dagger}$\,Zhao \& Ma (2025) analyse SHAP instability in an AutoML time
   series ensemble and propose a surrogate-model fix; they do not frame their
   analysis in terms of the Rashomon set.};

\end{tikzpicture}
